\chapter{The Harmonic Series}

\section{What is the harmonic series?}
The harmonic series is the endless sum
\[
1+\frac12+\frac13+\frac14+\frac15+\cdots
\]
It keeps adding reciprocals of positive whole numbers.

At first, the terms become small very quickly, so it is natural to wonder:
does this sum settle down to a fixed number?

\section{A quick guess from partial sums}
Let us add a few terms:
\[
1=1,\quad
1+\frac12=1.5,\quad
1+\frac12+\frac13\approx1.833,\quad
1+\frac12+\frac13+\frac14\approx2.083.
\]
It grows slowly, but it does keep growing.

\section{Proof of divergence: Comparison Test}
Now we prove that it never settles to a finite limit.

Start from
\[
1+\frac12+\frac13+\frac14+\frac15+\frac16+\frac17+\frac18+\cdots
\]
and group terms like this:
\[
1
\left(\frac12\right)
\left(\frac13+\frac14\right)
\left(\frac15+\frac16+\frac17+\frac18\right)
\cdots
\]

In each group, every term is at least as large as the last one in that group:
\begin{itemize}
  \item In \(\frac13+\frac14\), each term is at least \(\frac14\), so the group is at least
  \[
  \frac14+\frac14=\frac12.
  \]
  \item In \(\frac15+\frac16+\frac17+\frac18\), each term is at least \(\frac18\), so the group is at least
  \[
  \frac18+\frac18+\frac18+\frac18=\frac12.
  \]
\end{itemize}

The same pattern continues: each new group has twice as many terms as the previous
one, and each full group contributes at least \(\frac12\).

So the harmonic series is bigger than
\[
1+\frac12+\frac12+\frac12+\frac12+\cdots
\]
But this comparison sum clearly grows without bound. Therefore the harmonic
series also grows without bound.

\section{Infinite Prime Numbers: Alternativ Method}
Each integer \(n\) greater than 1 can be written as a product of primes:
\[
n=p_1^{a_1}p_2^{a_2}\cdots p_k^{a_k}
\]

\[
    1+\frac12+\frac13+\frac14+\frac15+\cdots
    = (1+\frac1{2}+\frac1{2^2}+...)(1+\frac1{3}+\frac1{3^2}+...)(1+\frac1{5}+\frac1{5^2}+...)...
\]

This identity becomes clear when you expand the product on the right: every
positive integer appears exactly once through its prime-factor form.

We know the sum on the left diverges. So the product on the right cannot
converge to a finite number either. That requires infinitely many primes
to appear: if only finitely many primes occurred, the product would reduce
to a finite product of convergent geometric series and would converge.

We can see here the unexpected connection between the harmonic series and the prime numbers.
That is the beauty of mathematics: different concepts can be connected in unexpected ways.
Dear reader, it will be thrilling to discover more such connections in your future studies.
Have fun!